\section*{On randomized reductions to $K^t$ (author Section 7)}
In this section, we show limitations on randomized nonadaptive reductions to $K^t$ for polynomially bounded $t$. We first need to define the notion of a hitting set generator.
\paragraph{Definition 33.} Given integers $n$ and $\ell<n$ and rational error parameter $\epsilon\in[0,1]$ with $\epsilon\ge1/2^n$, a hitting set generator $H_{n,\epsilon}$ for size $n$ with seed length $\ell$ and error $\epsilon$ is a function from $\ell(n)$ bits to $n$ bits, such that for any circuit $C_n$ of size $n$, if $C_n$ accepts at least an $\epsilon$ fraction of inputs of length $n$, we have that $\Pr_{z\in\{0,1\}^\ell}[C_n(H_{n,\epsilon}(z))=1]>0$. Given a function $\ell:\mathbb N\times[0,1]\to\mathbb N$, we say that a sequence $H_{n,\epsilon}$ of functions is a polynomial-time computable hitting set generator with seed length $\ell$ if for each integer $n$ and rational $\epsilon\in[0,1]$ such that $\epsilon\ge1/2^n$, $H_{n,\epsilon}$ is a hitting set generator for size $n$ with error $\epsilon$ and seed length $\ell(n,\epsilon)$, and moreover $H_{n,\epsilon}$ can be computed in time $\operatorname{poly}(n,\log(1/\epsilon))$ when given $n$ and $\lceil\log(1/\epsilon)\rceil$ in unary.
A key component of our proof is a construction of hitting set generators with optimal dependence on $\epsilon$ in the seed length [19].
\paragraph{Theorem 34 ([19]).} If E requires non-deterministic circuits of size $2^{\Omega(n)}$, then there is a polynomial-time computable hitting set generator $H$ with seed length $O(\log(n))+\log(1/\epsilon)$.
It is crucial for our results that the constant factor in front of the $\log(1/\epsilon)$ term is 1; if we could afford an arbitrary constant factor there, we would only need an assumption on hardness for deterministic circuits. We note that the result of [19] is stated for $\epsilon$ with a fixed dependence on $n$, but it is easy to see that their proof works also if $\log(1/\epsilon)$ is given as an independent parameter to the algorithm computing the hitting set generator.
We next use Theorem 34 to argue that under a standard derandomization assumptions, the Kolmogorov complexity of a typical sample from a polynomial-time samplable distributions does not differ by much from the time-bounded Kolmogorov complexity. We require that the time bound when measuring time-bounded Kolmogorov complexity is large enough compared to the time required for sampling.
\paragraph{Lemma 35.} Let $t':\mathbb N\to\mathbb N$ be a polynomially bounded function, and $\{D_m\}$ be a sequence of distributions samplable in time $t'(m)$, where for each $m\in\mathbb N$, $D_m$ is supported on $m$-bit strings. Suppose that E requires non-deterministic circuits of size $2^{\Omega(n)}$. Then there is $t:\mathbb N\to\mathbb N$ such that $t=\operatorname{poly}(t')$ and $\Pr_{x\sim D_m}[K^t(x)-K(x)=\omega(\log(m))]=1/m^{\omega(1)}$.
\begin{proof}Let $\{D_m\}$ be a sequence of distributions, where $D_m$ is supported on $m$-bit strings, sampled by an algorithm $A$ running in time $t'(m)$. Using the standard simulation of deterministic time by circuit size, $A$ can be simulated by a sequence of circuits $\{C'_m\}$, where $C'_m$ is of size at most $t'(m)^2$ for large enough $m$. Also, let $f=\omega(\log(m))$ be an arbitrary function.
Let $z$ be a fixed string of length at most $n=2t'(m)^2$ that is assigned probability $\epsilon$ by $D_m$. We argue that the polynomial-time bounded Kolmogorov complexity of $z$ is at most $\log(1/\epsilon)+O(\log(n))$, under the assumption that E required non-deterministic circuits of size $2^{\Omega(n)}$. Indeed, under the given derandomization assumption, there is hitting set generator $H=\{H_{n,\epsilon}\}$ with seed length $O(\log(n))+\log(1/\epsilon))$ computable in polynomial time as per Definition 33. Let $B$ be an algorithm that computes $H$ in polynomial time.
We show that there must be at least one output $r$ of $H_{n,\epsilon}$ such that $C'_m$ outputs $z$ on input $r$. Indeed, consider a circuit $C_n$ that on input $r$ of length at most $n$, runs $C'_m(r')$ and accepts iff the answer is $z$. Since $z$ has probability $\epsilon$ according to $D_m$ and $C'_m$ samples from $D_m$, we have that $C_n$ accepts with probability at least $\epsilon$ on a uniformly random input. Since $H_{n,\epsilon}$ is a hitting set generator, we have that at least one output $r$ of the hitting set generator is such that $C'_m(r)=z$.
Let $y$ be the seed of length $O(\log(n))+\log(1/\epsilon)$ for which $H_{n,\epsilon}(y)=r$. We describe $z$ using the code of $B$, the string $y$ and descriptions of $n$ and $\lceil\log(1/\epsilon)\rceil$ in binary. This composite description is of size at most $O(\log(n))+\log(1/\epsilon)$, and $z$ can be computed in time $t=\operatorname{poly}(n,\log(1/\epsilon))$ from this description by running the algorithm $B$ with parameters $n$ and $\log(1/\epsilon)$ and input $y$ to produce a string $r$, and then running $A$ on $r$ to produce $z$. Hence $K^t(z)=O(\log(n))+\log(1/\epsilon)$. Let $C$ be an explicit constant such that $K^t(z)\le C\log(m)+\log(1/\epsilon)$ for large enough $n$.
In the next part of our argument, we show that the cumulative probability according to $D_m$ of strings with additive gap at least $f(m)$ between Kolmogorov complexity and $t$-time bounded Kolmogorov complexity is $1/m^{\omega(1)}$, establishing our result.
Partition all binary strings of length $m$ into $2m$ groups as follows. The group $S_i$, $i\in[2m-1]$ contains all strings of length $m$ sampled with probability at most $1/2^{i-1}$ and greater than $1/2^i$ from $D_m$. The group $S_m$ contains all strings of length $m$ sampled with probability at most $1/2^{2m-1}$. We will obtain our upper bound on the likelihood of a large gap between $K$ and $Kt$ complexity by using a union bound on the groups $S_i$.
Consider any group $S_i$, $i\in[2m-1]$. Let $g(m)=f(m)-C\log(m)$. Since $f(m)=\omega(\log(m))$, so is $g(m)$. Call a string $x$ of length $m$ $i$-bad if $x\in S_i$ and $K(x)\le i-g(m)$. By the upper bound on Komogorov complexity, we have that there are at most $2^{i-g(m)}=2^i/m^{\omega(1)}$ $i$-bad strings. Since each $i$-bad string is in $S_i$, it has probability at most $1/2^{i-1}$ of being sampled according to $D_m$, and therefore the cumulative probability of $i$-bad strings according to $D_m$ is at most $1/2^{i-1}\cdot2^i/m^{\omega(1)}=1/m^{\omega(1)}$.
By the argument earlier in the proof, each string in $S_i$ has $K^t$ complexity at most $C\log(m)+\log(1/2^i)=C\log(m)+i$. Hence any string $x\in S_i$ such that $K^t(x)-K(x)\ge f(m)$ is $i$-bad, and the cumulative probability of all such strings according to $D_m$ is $1/m^{\omega(1)}$. By a union bound over $i\in[2m-1]$, we have that the cumulative probability of all strings with gap at least $f(m)$ between Kolmogorov complexity and $t$-time bounded Kolmogorov complexity and that occur with probability greater than $1/2^{2m}$ in $D_m$ is $1/m^{\omega(1)}$. Since there are only $2^m$ $m$-bit strings, the cumulative probability according to $D_m$ of strings from $S_{2m}$ is at most $2^{-m}=1/m^{\omega(1)}$. Hence, by another union bound, $\Pr_{x\sim D_m}[K^t(x)-K(x)=\omega(\log(m))]=1/m^{\omega(1)}$, as desired.\end{proof}
We note that a weaker version of Lemma 35, where the hypothesis is that E requires $\Sigma_2^p$-oracle circuits of truly exponential size, is implicit in work of Antunes and Fortnow [16]. We require a couple of other well-known results. The first is a proposition shown using a straightforward padding argument.
\paragraph{Proposition 36.} If every tally language in $\mathrm{NTIME}(n)$ is in coNP, then $\mathrm{NE}=\mathrm{coNE}$.
The second is a standard derandomization result.
\paragraph{Theorem 37 ([38]).} If there is an $\epsilon>0$ such that E does not have non-deterministic circuits of size $2^{\epsilon n}$, then $\mathrm{AM}=\mathrm{NP}$.
Now we have the tools to state and prove the main result of this section, which also appears as Theorem 6 in the Introduction.
\paragraph{Proof of Theorem 38.} Suppose there is a randomized time $t'$ non-adaptive reduction $R$ from SAT to $O(\log(n))$-additively approximating $K^t$ complexity with the properties specified in the statement of the Theorem, where $t=\operatorname{poly}(t')$ and the precise polynomial dependence is to be specified later. We use Lemma 35 to show that under the assumption that E does not have non-deterministic circuits of size $2^{\epsilon n}$ for some $\epsilon>0$, $R$ also reduces any tally language in $\mathrm{NTIME}(n)$ to $\omega(\log(m))$-additively approximating $K$ complexity, and then apply Theorem 29.
Indeed, let $L\in\mathrm{NTIME}(n)$ be a tally language. By the Cook--Levin theorem, there is a m-reduction from $L$ to SAT running in time $n\operatorname{polylog}(n)$. This reduction is polynomially honest and has fixed query length. By composing this reduction with the reduction $R$ in the assumption, we get that there is a randomized time $t''$ non-adaptive reduction $R'$ from $L$ to $\omega(\log(m))$-additively approximating $K^t$ complexity, where the reduction is polynomially honest and has fixed query length, and where $t''=t(n\operatorname{polylog}(n))$ is polynomially bounded.
Now we apply Lemma 35 to the sequence of distributions $D_m$ sampled by running the randomized reduction $R'$ on input $1^m$ and outputting a random query.\footnote{Technically speaking, it might be unclear if Lemma 35 applies here, as the sampled strings do not in general have length $m$, but the proof of Lemma 35 continues to work as long as the sampled strings have length $m^{\Omega(1)}$.} Applying Lemma 35 and a union bound, we have that with probability $1-1/m^{\Omega(1)}$ over the randomness $r$ of $R'$, all queries $z$ asked on randomness $r$ have $|K(z)-K^t(z)|=O(\log(m))$, and since the reduction has inverse polynomial advantage, replacing an oracle that $\omega(\log(m))$-additively approximates $K^t$ complexity with one that $\omega(\log(m))$-additively approximates $K$ complexity will preserve the answer of the reduction for each string $1^m$. Hence we have that the hypothesis of Theorem 29 is satisfied, and we get that $L\in\mathrm{AM}\cap\mathrm{coAM}$.
Using the assumption that E does not have non-deterministic circuits of size $2^{\epsilon n}$ once more and applying Theorem 37, we have that $L\in\mathrm{NP}\cap\mathrm{coNP}$. Since this is the case for every $L\in\mathrm{NTIME}(n)$, we can use Proposition 36 and conclude that $\mathrm{NE}=\mathrm{coNE}$, as desired.\qed
